% figs_part1.tex —— 由 mit/make_figs_part1.py 產生，請勿手動修改。
% 重新產生：python mit/make_figs_part1.py（在 notes/ 裡）
\newcommand{\FigMiCauchy}{%
\begin{tikzpicture}
\node[font=\small,text=gray!70!black] at (0.625,0) {$\times$};
\node[font=\small,text=cRight] at (2.06,0) {$C_0$};
\node[font=\small,text=cRight] at (3.68,0) {$C_1x$};
\node[font=\small,text=cRight] at (5.3,0) {$C_2x^2$};
\node[font=\small,text=cRight] at (6.92,0) {$C_3x^3$};
\node[font=\small,text=cLeft] at (0.625,-0.78) {$C_0$};
\node[font=\small,text=cLeft] at (0.625,-1.56) {$C_1x$};
\node[font=\small,text=cLeft] at (0.625,-2.34) {$C_2x^2$};
\node[font=\small,text=cLeft] at (0.625,-3.12) {$C_3x^3$};
\fill[cA!9] (1.25,-1.17) rectangle (2.87,-0.39);
\node[font=\small] at (2.06,-0.65) {$\textcolor{cLeft}{C_0}\textcolor{cRight}{C_0}$};
\node[font=\scriptsize,text=gray!70!black] at (2.06,-0.97) {$=1$};
\fill[green!50!black!11] (2.87,-1.17) rectangle (4.49,-0.39);
\node[font=\small] at (3.68,-0.65) {$\textcolor{cLeft}{C_0}\textcolor{cRight}{C_1}x$};
\node[font=\scriptsize,text=gray!70!black] at (3.68,-0.97) {$=x$};
\fill[yellow!45] (4.49,-1.17) rectangle (6.11,-0.39);
\node[font=\small] at (5.3,-0.65) {$\textcolor{cLeft}{C_0}\textcolor{cRight}{C_2}x^2$};
\node[font=\scriptsize,text=gray!70!black] at (5.3,-0.97) {$=2x^2$};
\fill[orange!20] (6.11,-1.17) rectangle (7.73,-0.39);
\node[font=\small] at (6.92,-0.65) {$\textcolor{cLeft}{C_0}\textcolor{cRight}{C_3}x^3$};
\node[font=\scriptsize,text=gray!70!black] at (6.92,-0.97) {$=5x^3$};
\fill[green!50!black!11] (1.25,-1.95) rectangle (2.87,-1.17);
\node[font=\small] at (2.06,-1.43) {$\textcolor{cLeft}{C_1}\textcolor{cRight}{C_0}x$};
\node[font=\scriptsize,text=gray!70!black] at (2.06,-1.75) {$=x$};
\fill[yellow!45] (2.87,-1.95) rectangle (4.49,-1.17);
\node[font=\small] at (3.68,-1.43) {$\textcolor{cLeft}{C_1}\textcolor{cRight}{C_1}x^2$};
\node[font=\scriptsize,text=gray!70!black] at (3.68,-1.75) {$=x^2$};
\fill[orange!20] (4.49,-1.95) rectangle (6.11,-1.17);
\node[font=\small] at (5.3,-1.43) {$\textcolor{cLeft}{C_1}\textcolor{cRight}{C_2}x^3$};
\node[font=\scriptsize,text=gray!70!black] at (5.3,-1.75) {$=2x^3$};
\node[font=\small,text=gray!55] at (6.92,-1.43) {$C_1C_3x^4$};
\node[font=\scriptsize,text=gray!55] at (6.92,-1.75) {$=5x^4$};
\fill[yellow!45] (1.25,-2.73) rectangle (2.87,-1.95);
\node[font=\small] at (2.06,-2.21) {$\textcolor{cLeft}{C_2}\textcolor{cRight}{C_0}x^2$};
\node[font=\scriptsize,text=gray!70!black] at (2.06,-2.53) {$=2x^2$};
\fill[orange!20] (2.87,-2.73) rectangle (4.49,-1.95);
\node[font=\small] at (3.68,-2.21) {$\textcolor{cLeft}{C_2}\textcolor{cRight}{C_1}x^3$};
\node[font=\scriptsize,text=gray!70!black] at (3.68,-2.53) {$=2x^3$};
\node[font=\small,text=gray!55] at (5.3,-2.21) {$C_2C_2x^4$};
\node[font=\scriptsize,text=gray!55] at (5.3,-2.53) {$=4x^4$};
\node[font=\small,text=gray!55] at (6.92,-2.21) {$C_2C_3x^5$};
\node[font=\scriptsize,text=gray!55] at (6.92,-2.53) {$=10x^5$};
\fill[orange!20] (1.25,-3.51) rectangle (2.87,-2.73);
\node[font=\small] at (2.06,-2.99) {$\textcolor{cLeft}{C_3}\textcolor{cRight}{C_0}x^3$};
\node[font=\scriptsize,text=gray!70!black] at (2.06,-3.31) {$=5x^3$};
\node[font=\small,text=gray!55] at (3.68,-2.99) {$C_3C_1x^4$};
\node[font=\scriptsize,text=gray!55] at (3.68,-3.31) {$=5x^4$};
\node[font=\small,text=gray!55] at (5.3,-2.99) {$C_3C_2x^5$};
\node[font=\scriptsize,text=gray!55] at (5.3,-3.31) {$=10x^5$};
\node[font=\small,text=gray!55] at (6.92,-2.99) {$C_3C_3x^6$};
\node[font=\scriptsize,text=gray!55] at (6.92,-3.31) {$=25x^6$};
\draw[gray!55,line width=0.4pt] (1.25,-0.39) -- (7.73,-0.39);
\draw[gray!55,line width=0.4pt] (1.25,-1.17) -- (7.73,-1.17);
\draw[gray!55,line width=0.4pt] (1.25,-1.95) -- (7.73,-1.95);
\draw[gray!55,line width=0.4pt] (1.25,-2.73) -- (7.73,-2.73);
\draw[gray!55,line width=0.4pt] (1.25,-3.51) -- (7.73,-3.51);
\draw[gray!55,line width=0.4pt] (1.25,-0.39) -- (1.25,-3.51);
\draw[gray!55,line width=0.4pt] (2.87,-0.39) -- (2.87,-3.51);
\draw[gray!55,line width=0.4pt] (4.49,-0.39) -- (4.49,-3.51);
\draw[gray!55,line width=0.4pt] (6.11,-0.39) -- (6.11,-3.51);
\draw[gray!55,line width=0.4pt] (7.73,-0.39) -- (7.73,-3.51);
\fill[cA!9] (8.48,-0.95) rectangle (8.82,-0.61);
\draw[gray!55,line width=0.3pt] (8.48,-0.95) rectangle (8.82,-0.61);
\node[anchor=west,font=\small] at (8.98,-0.78) {$[x^0]\,C^2=1$};
\fill[green!50!black!11] (8.48,-1.73) rectangle (8.82,-1.39);
\draw[gray!55,line width=0.3pt] (8.48,-1.73) rectangle (8.82,-1.39);
\node[anchor=west,font=\small] at (8.98,-1.56) {$[x]\,C^2=1+1=2$};
\fill[yellow!45] (8.48,-2.51) rectangle (8.82,-2.17);
\draw[gray!55,line width=0.3pt] (8.48,-2.51) rectangle (8.82,-2.17);
\node[anchor=west,font=\small] at (8.98,-2.34) {$[x^2]\,C^2=2+1+2=5$};
\fill[orange!20] (8.48,-3.29) rectangle (8.82,-2.95);
\draw[gray!55,line width=0.3pt] (8.48,-3.29) rectangle (8.82,-2.95);
\node[anchor=west,font=\small] at (8.98,-3.12) {$[x^3]\,C^2=5+2+2+5=14$};
\node[anchor=west,font=\footnotesize,text=gray!70!black] at (8.48,0) {同色的格子加起來：};
\end{tikzpicture}}
\newcommand{\FigMiiBranches}{%
\begin{tikzpicture}
\draw[-{Stealth[length=4pt]},gray!70!black,line width=0.5pt] (0,0) -- (4.2,0) node[right,font=\small] {$x$};
\draw[-{Stealth[length=4pt]},gray!70!black,line width=0.5pt] (0,0) -- (0,2.85);
\draw[gray!45,dashed,line width=0.4pt] (3.75,0) -- (3.75,1);
\node[below,font=\small] at (3.75,0) {$\frac14$};
\node[below left,font=\small] at (0,0) {$0$};
\draw[gray!70!black,line width=0.4pt] (-0.07,0.5) -- (0.07,0.5);
\node[left,font=\scriptsize] at (-0.07,0.5) {$1$};
\draw[gray!70!black,line width=0.4pt] (-0.07,1) -- (0.07,1);
\node[left,font=\scriptsize] at (-0.07,1) {$2$};
\draw[gray!70!black,line width=0.4pt] (-0.07,1.5) -- (0.07,1.5);
\node[left,font=\scriptsize] at (-0.07,1.5) {$3$};
\draw[gray!70!black,line width=0.4pt] (-0.07,2) -- (0.07,2);
\node[left,font=\scriptsize] at (-0.07,2) {$4$};
\draw[gray!70!black,line width=0.4pt] (-0.07,2.5) -- (0.07,2.5);
\node[left,font=\scriptsize] at (-0.07,2.5) {$5$};
\draw[cA,line width=1.3pt] (0,0.5) -- (0.037,0.501) -- (0.075,0.503) -- (0.112,0.504) -- (0.15,0.505) -- (0.188,0.506) -- (0.225,0.508) -- (0.263,0.509) -- (0.3,0.51) -- (0.337,0.512) -- (0.375,0.513) -- (0.412,0.515) -- (0.45,0.516) -- (0.488,0.517) -- (0.525,0.519) -- (0.562,0.52) -- (0.6,0.522) -- (0.638,0.523) -- (0.675,0.525) -- (0.713,0.526) -- (0.75,0.528) -- (0.787,0.529) -- (0.825,0.531) -- (0.863,0.533) -- (0.9,0.534) -- (0.938,0.536) -- (0.975,0.538) -- (1.013,0.539) -- (1.05,0.541) -- (1.087,0.543) -- (1.125,0.544) -- (1.163,0.546) -- (1.2,0.548) -- (1.238,0.55) -- (1.275,0.552) -- (1.313,0.554) -- (1.35,0.556) -- (1.387,0.557) -- (1.425,0.559) -- (1.463,0.561) -- (1.5,0.564) -- (1.538,0.566) -- (1.575,0.568) -- (1.613,0.57) -- (1.65,0.572) -- (1.688,0.574) -- (1.725,0.576) -- (1.763,0.579) -- (1.8,0.581) -- (1.837,0.583) -- (1.875,0.586) -- (1.913,0.588) -- (1.95,0.591) -- (1.988,0.593) -- (2.025,0.596) -- (2.062,0.599) -- (2.1,0.601) -- (2.138,0.604) -- (2.175,0.607) -- (2.212,0.61) -- (2.25,0.613) -- (2.288,0.616) -- (2.325,0.619) -- (2.362,0.622) -- (2.4,0.625) -- (2.438,0.628) -- (2.475,0.632) -- (2.513,0.635) -- (2.55,0.639) -- (2.588,0.642) -- (2.625,0.646) -- (2.662,0.65) -- (2.7,0.654) -- (2.737,0.658) -- (2.775,0.662) -- (2.812,0.667) -- (2.85,0.671) -- (2.888,0.676) -- (2.925,0.681) -- (2.963,0.686) -- (3,0.691) -- (3.038,0.696) -- (3.075,0.702) -- (3.113,0.708) -- (3.15,0.714) -- (3.188,0.721) -- (3.225,0.728) -- (3.263,0.735) -- (3.3,0.743) -- (3.337,0.751) -- (3.375,0.76) -- (3.413,0.769) -- (3.45,0.78) -- (3.488,0.791) -- (3.525,0.803) -- (3.563,0.817) -- (3.6,0.833) -- (3.637,0.852) -- (3.675,0.876) -- (3.712,0.909) -- (3.75,1);
\draw[cBad,line width=1.3pt] (2.4,2.5) -- (2.415,2.479) -- (2.43,2.459) -- (2.445,2.439) -- (2.46,2.418) -- (2.475,2.399) -- (2.49,2.379) -- (2.505,2.36) -- (2.52,2.34) -- (2.535,2.321) -- (2.55,2.302) -- (2.565,2.284) -- (2.58,2.265) -- (2.595,2.247) -- (2.61,2.229) -- (2.625,2.211) -- (2.64,2.193) -- (2.655,2.176) -- (2.67,2.158) -- (2.685,2.141) -- (2.7,2.124) -- (2.715,2.107) -- (2.73,2.09) -- (2.745,2.073) -- (2.76,2.057) -- (2.775,2.04) -- (2.79,2.024) -- (2.805,2.008) -- (2.82,1.992) -- (2.835,1.976) -- (2.85,1.96) -- (2.865,1.945) -- (2.88,1.929) -- (2.895,1.914) -- (2.91,1.899) -- (2.925,1.883) -- (2.94,1.868) -- (2.955,1.853) -- (2.97,1.838) -- (2.985,1.824) -- (3,1.809) -- (3.015,1.794) -- (3.03,1.78) -- (3.045,1.766) -- (3.06,1.751) -- (3.075,1.737) -- (3.09,1.723) -- (3.105,1.709) -- (3.12,1.695) -- (3.135,1.681) -- (3.15,1.667) -- (3.165,1.653) -- (3.18,1.639) -- (3.195,1.625) -- (3.21,1.612) -- (3.225,1.598) -- (3.24,1.584) -- (3.255,1.571) -- (3.27,1.557) -- (3.285,1.544) -- (3.3,1.53) -- (3.315,1.517) -- (3.33,1.503) -- (3.345,1.489) -- (3.36,1.476) -- (3.375,1.462) -- (3.39,1.449) -- (3.405,1.435) -- (3.42,1.422) -- (3.435,1.408) -- (3.45,1.394) -- (3.465,1.381) -- (3.48,1.367) -- (3.495,1.353) -- (3.51,1.339) -- (3.525,1.324) -- (3.54,1.31) -- (3.555,1.295) -- (3.57,1.281) -- (3.585,1.265) -- (3.6,1.25) -- (3.615,1.234) -- (3.63,1.218) -- (3.645,1.201) -- (3.66,1.183) -- (3.675,1.165) -- (3.69,1.145) -- (3.705,1.123) -- (3.72,1.098) -- (3.735,1.068) -- (3.75,1);
\draw[-{Stealth[length=5pt]},cBad,line width=1.3pt] (2.4,2.5) -- (2.22,2.768);
\fill[cA] (0,0.5) circle (0.07);
\fill[black] (3.75,1) circle (0.05);
\node[anchor=north west,font=\scriptsize,text=cA] at (0.04,0.46) {$C(0)=C_0=1$};
\node[right,font=\scriptsize,text=cBad] at (2.475,2.375) {取 $+$：$x\to0$ 時 $\to\infty$};
\node[above,font=\scriptsize,text=cA] at (1.95,0.641) {取 $-$};
\node[right,font=\scriptsize] at (3.8,1.12) {$(\frac14,2)$};
\end{tikzpicture}}
\newcommand{\FigMiiiWalk}{%
\begin{tikzpicture}
\draw[gray!22,line width=0.3pt] (0,0.42) -- (3.74,0.42);
\draw[gray!22,line width=0.3pt] (0,0.84) -- (3.74,0.84);
\draw[gray!60,line width=0.5pt] (0,0) -- (3.74,0);
\draw[cBad,dashed,line width=0.6pt] (0,-0.42) -- (3.74,-0.42) node[right,font=\scriptsize,text=cBad] {懸崖};
\draw[-{Stealth[length=4pt]},gray!70!black,line width=0.5pt] (0,-0.672) -- (0,1.26) node[above,font=\scriptsize] {位置};
\node[font=\scriptsize,text=gray!70!black] at (1.82,-1.155) {時間（步數）$\to$};
\node[left,font=\scriptsize,text=cBad] at (-0.02,-0.42) {$-1$};
\node[left,font=\scriptsize,text=black] at (-0.02,0) {$0$};
\node[left,font=\scriptsize,text=black] at (-0.02,0.42) {$1$};
\node[left,font=\scriptsize,text=black] at (-0.02,0.84) {$2$};
\draw[cA,line width=1.5pt,line cap=round] (0,0) -- (0.52,0.42);
\draw[cA,line width=1.5pt,line cap=round] (0.52,0.42) -- (1.04,0.84);
\draw[cA,line width=1.5pt,line cap=round] (1.04,0.84) -- (1.56,0.42);
\draw[cA,line width=1.5pt,line cap=round] (1.56,0.42) -- (2.08,0.84);
\draw[cA,line width=1.5pt,line cap=round] (2.08,0.84) -- (2.6,0.42);
\draw[cA,line width=1.5pt,line cap=round] (2.6,0.42) -- (3.12,0);
\draw[cBad,line width=1.5pt,line cap=round] (3.12,0) -- (3.64,-0.42);
\fill[black!70] (0,0) circle (0.035);
\fill[black!70] (0.52,0.42) circle (0.035);
\fill[black!70] (1.04,0.84) circle (0.035);
\fill[black!70] (1.56,0.42) circle (0.035);
\fill[black!70] (2.08,0.84) circle (0.035);
\fill[black!70] (2.6,0.42) circle (0.035);
\fill[black!70] (3.12,0) circle (0.035);
\fill[black!70] (3.64,-0.42) circle (0.035);
\fill[cBad] (3.64,-0.42) circle (0.07);
\node[font=\scriptsize,text=cA] at (0.26,-0.777) {$\mathrm{U}$};
\node[font=\scriptsize,text=cA] at (0.78,-0.777) {$\mathrm{U}$};
\node[font=\scriptsize,text=cA] at (1.3,-0.777) {$\mathrm{D}$};
\node[font=\scriptsize,text=cA] at (1.82,-0.777) {$\mathrm{U}$};
\node[font=\scriptsize,text=cA] at (2.34,-0.777) {$\mathrm{D}$};
\node[font=\scriptsize,text=cA] at (2.86,-0.777) {$\mathrm{D}$};
\node[font=\scriptsize,text=cBad] at (3.38,-0.777) {$\mathrm{D}$};
\node[font=\small] at (5.4,1.05) {$n=0$};
\draw[gray!55,line width=0.3pt] (5.32,0.4) -- (5.48,0.4);
\draw[cBad!70,dashed,line width=0.3pt] (5.32,0.24) -- (5.48,0.24);
\draw[cBad,line width=1.1pt,line cap=round] (5.32,0.4) -- (5.48,0.24);
\node[font=\small] at (5.4,-1) {$1$ 條};
\node[font=\small] at (5.4,-1.6) {$\frac{1}{2}$};
\node[font=\small] at (6.65,1.05) {$n=1$};
\draw[gray!55,line width=0.3pt] (6.41,0.4) -- (6.89,0.4);
\draw[cBad!70,dashed,line width=0.3pt] (6.41,0.24) -- (6.89,0.24);
\draw[cA,line width=1.1pt,line cap=round] (6.41,0.4) -- (6.57,0.56);
\draw[cA,line width=1.1pt,line cap=round] (6.57,0.56) -- (6.73,0.4);
\draw[cBad,line width=1.1pt,line cap=round] (6.73,0.4) -- (6.89,0.24);
\node[font=\small] at (6.65,-1) {$1$ 條};
\node[font=\small] at (6.65,-1.6) {$\frac{1}{8}$};
\node[font=\small] at (8.55,1.05) {$n=2$};
\draw[gray!55,line width=0.3pt] (7.64,0.4) -- (8.44,0.4);
\draw[cBad!70,dashed,line width=0.3pt] (7.64,0.24) -- (8.44,0.24);
\draw[cA,line width=1.1pt,line cap=round] (7.64,0.4) -- (7.8,0.56);
\draw[cA,line width=1.1pt,line cap=round] (7.8,0.56) -- (7.96,0.4);
\draw[cA,line width=1.1pt,line cap=round] (7.96,0.4) -- (8.12,0.56);
\draw[cA,line width=1.1pt,line cap=round] (8.12,0.56) -- (8.28,0.4);
\draw[cBad,line width=1.1pt,line cap=round] (8.28,0.4) -- (8.44,0.24);
\draw[gray!55,line width=0.3pt] (8.66,0.4) -- (9.46,0.4);
\draw[cBad!70,dashed,line width=0.3pt] (8.66,0.24) -- (9.46,0.24);
\draw[cA,line width=1.1pt,line cap=round] (8.66,0.4) -- (8.82,0.56);
\draw[cA,line width=1.1pt,line cap=round] (8.82,0.56) -- (8.98,0.72);
\draw[cA,line width=1.1pt,line cap=round] (8.98,0.72) -- (9.14,0.56);
\draw[cA,line width=1.1pt,line cap=round] (9.14,0.56) -- (9.3,0.4);
\draw[cBad,line width=1.1pt,line cap=round] (9.3,0.4) -- (9.46,0.24);
\node[font=\small] at (8.55,-1) {$2$ 條};
\node[font=\small] at (8.55,-1.6) {$\frac{2}{32}$};
\node[font=\small] at (12.11,1.05) {$n=3$};
\draw[gray!55,line width=0.3pt] (10.21,0.4) -- (11.33,0.4);
\draw[cBad!70,dashed,line width=0.3pt] (10.21,0.24) -- (11.33,0.24);
\draw[cA,line width=1.1pt,line cap=round] (10.21,0.4) -- (10.37,0.56);
\draw[cA,line width=1.1pt,line cap=round] (10.37,0.56) -- (10.53,0.4);
\draw[cA,line width=1.1pt,line cap=round] (10.53,0.4) -- (10.69,0.56);
\draw[cA,line width=1.1pt,line cap=round] (10.69,0.56) -- (10.85,0.4);
\draw[cA,line width=1.1pt,line cap=round] (10.85,0.4) -- (11.01,0.56);
\draw[cA,line width=1.1pt,line cap=round] (11.01,0.56) -- (11.17,0.4);
\draw[cBad,line width=1.1pt,line cap=round] (11.17,0.4) -- (11.33,0.24);
\draw[gray!55,line width=0.3pt] (11.55,0.4) -- (12.67,0.4);
\draw[cBad!70,dashed,line width=0.3pt] (11.55,0.24) -- (12.67,0.24);
\draw[cA,line width=1.1pt,line cap=round] (11.55,0.4) -- (11.71,0.56);
\draw[cA,line width=1.1pt,line cap=round] (11.71,0.56) -- (11.87,0.4);
\draw[cA,line width=1.1pt,line cap=round] (11.87,0.4) -- (12.03,0.56);
\draw[cA,line width=1.1pt,line cap=round] (12.03,0.56) -- (12.19,0.72);
\draw[cA,line width=1.1pt,line cap=round] (12.19,0.72) -- (12.35,0.56);
\draw[cA,line width=1.1pt,line cap=round] (12.35,0.56) -- (12.51,0.4);
\draw[cBad,line width=1.1pt,line cap=round] (12.51,0.4) -- (12.67,0.24);
\draw[gray!55,line width=0.3pt] (12.89,0.4) -- (14.01,0.4);
\draw[cBad!70,dashed,line width=0.3pt] (12.89,0.24) -- (14.01,0.24);
\draw[cA,line width=1.1pt,line cap=round] (12.89,0.4) -- (13.05,0.56);
\draw[cA,line width=1.1pt,line cap=round] (13.05,0.56) -- (13.21,0.72);
\draw[cA,line width=1.1pt,line cap=round] (13.21,0.72) -- (13.37,0.56);
\draw[cA,line width=1.1pt,line cap=round] (13.37,0.56) -- (13.53,0.4);
\draw[cA,line width=1.1pt,line cap=round] (13.53,0.4) -- (13.69,0.56);
\draw[cA,line width=1.1pt,line cap=round] (13.69,0.56) -- (13.85,0.4);
\draw[cBad,line width=1.1pt,line cap=round] (13.85,0.4) -- (14.01,0.24);
\draw[cA!60,rounded corners=1.5pt,line width=0.5pt] (10.8,-0.74) rectangle (12.08,-0.1);
\draw[gray!55,line width=0.3pt] (10.88,-0.5) -- (12,-0.5);
\draw[cBad!70,dashed,line width=0.3pt] (10.88,-0.66) -- (12,-0.66);
\draw[cA,line width=1.1pt,line cap=round] (10.88,-0.5) -- (11.04,-0.34);
\draw[cA,line width=1.1pt,line cap=round] (11.04,-0.34) -- (11.2,-0.18);
\draw[cA,line width=1.1pt,line cap=round] (11.2,-0.18) -- (11.36,-0.34);
\draw[cA,line width=1.1pt,line cap=round] (11.36,-0.34) -- (11.52,-0.18);
\draw[cA,line width=1.1pt,line cap=round] (11.52,-0.18) -- (11.68,-0.34);
\draw[cA,line width=1.1pt,line cap=round] (11.68,-0.34) -- (11.84,-0.5);
\draw[cBad,line width=1.1pt,line cap=round] (11.84,-0.5) -- (12,-0.66);
\draw[gray!55,line width=0.3pt] (12.22,-0.5) -- (13.34,-0.5);
\draw[cBad!70,dashed,line width=0.3pt] (12.22,-0.66) -- (13.34,-0.66);
\draw[cA,line width=1.1pt,line cap=round] (12.22,-0.5) -- (12.38,-0.34);
\draw[cA,line width=1.1pt,line cap=round] (12.38,-0.34) -- (12.54,-0.18);
\draw[cA,line width=1.1pt,line cap=round] (12.54,-0.18) -- (12.7,-0.02);
\draw[cA,line width=1.1pt,line cap=round] (12.7,-0.02) -- (12.86,-0.18);
\draw[cA,line width=1.1pt,line cap=round] (12.86,-0.18) -- (13.02,-0.34);
\draw[cA,line width=1.1pt,line cap=round] (13.02,-0.34) -- (13.18,-0.5);
\draw[cBad,line width=1.1pt,line cap=round] (13.18,-0.5) -- (13.34,-0.66);
\node[font=\small] at (12.11,-1) {$5$ 條};
\node[font=\small] at (12.11,-1.6) {$\frac{5}{128}$};
\node[anchor=east,font=\scriptsize,text=gray!70!black] at (5,0.4) {走法};
\node[anchor=east,font=\scriptsize,text=gray!70!black] at (5,-1) {條數};
\node[anchor=east,font=\scriptsize,text=gray!70!black] at (5,-1.6) {機率};
\end{tikzpicture}}
\newcommand{\FigMivRotations}{%
\begin{tikzpicture}
\draw[gray!60,dashed,line width=0.4pt] (-0.08,0) -- (1.55,0);
\draw[cA,line width=1.2pt,line cap=round] (0,0) -- (0.21,-0.21);
\draw[gray!75!black,line width=1.2pt,line cap=round] (0.21,-0.21) -- (0.42,0);
\draw[cA,line width=1.2pt,line cap=round] (0.42,0) -- (0.63,-0.21);
\draw[cA,line width=1.2pt,line cap=round] (0.63,-0.21) -- (0.84,-0.42);
\draw[gray!75!black,line width=1.2pt,line cap=round] (0.84,-0.42) -- (1.05,-0.21);
\draw[gray!75!black,line width=1.2pt,line cap=round] (1.05,-0.21) -- (1.26,0);
\draw[gray!75!black,line width=1.2pt,line cap=round] (1.26,0) -- (1.47,0.21);
\fill[black!70] (0,0) circle (0.028);
\fill[black!70] (0.21,-0.21) circle (0.028);
\fill[black!70] (0.42,0) circle (0.028);
\fill[black!70] (0.63,-0.21) circle (0.028);
\fill[black!70] (0.84,-0.42) circle (0.028);
\fill[black!70] (1.05,-0.21) circle (0.028);
\fill[black!70] (1.26,0) circle (0.028);
\fill[black!70] (1.47,0.21) circle (0.028);
\fill[cBad] (0.21,-0.21) circle (0.075);
\draw[cTop,line width=0.9pt] (0.84,-0.42) circle (0.11);
\node[font=\small] at (0.735,-0.865) {$j=0$\ \textcolor{cBad}{$\times$}};
\draw[gray!60,dashed,line width=0.4pt] (2.22,0) -- (3.85,0);
\draw[gray!75!black,line width=1.2pt,line cap=round] (2.3,0) -- (2.51,0.21);
\draw[cA,line width=1.2pt,line cap=round] (2.51,0.21) -- (2.72,0);
\draw[cA,line width=1.2pt,line cap=round] (2.72,0) -- (2.93,-0.21);
\draw[gray!75!black,line width=1.2pt,line cap=round] (2.93,-0.21) -- (3.14,0);
\draw[gray!75!black,line width=1.2pt,line cap=round] (3.14,0) -- (3.35,0.21);
\draw[gray!75!black,line width=1.2pt,line cap=round] (3.35,0.21) -- (3.56,0.42);
\draw[cA,line width=1.2pt,line cap=round] (3.56,0.42) -- (3.77,0.21);
\fill[black!70] (2.3,0) circle (0.028);
\fill[black!70] (2.51,0.21) circle (0.028);
\fill[black!70] (2.72,0) circle (0.028);
\fill[black!70] (2.93,-0.21) circle (0.028);
\fill[black!70] (3.14,0) circle (0.028);
\fill[black!70] (3.35,0.21) circle (0.028);
\fill[black!70] (3.56,0.42) circle (0.028);
\fill[black!70] (3.77,0.21) circle (0.028);
\fill[cBad] (2.72,0) circle (0.075);
\node[font=\small] at (3.035,-0.865) {$j=1$\ \textcolor{cBad}{$\times$}};
\draw[gray!60,dashed,line width=0.4pt] (4.52,0) -- (6.15,0);
\draw[cA,line width=1.2pt,line cap=round] (4.6,0) -- (4.81,-0.21);
\draw[cA,line width=1.2pt,line cap=round] (4.81,-0.21) -- (5.02,-0.42);
\draw[gray!75!black,line width=1.2pt,line cap=round] (5.02,-0.42) -- (5.23,-0.21);
\draw[gray!75!black,line width=1.2pt,line cap=round] (5.23,-0.21) -- (5.44,0);
\draw[gray!75!black,line width=1.2pt,line cap=round] (5.44,0) -- (5.65,0.21);
\draw[cA,line width=1.2pt,line cap=round] (5.65,0.21) -- (5.86,0);
\draw[gray!75!black,line width=1.2pt,line cap=round] (5.86,0) -- (6.07,0.21);
\fill[black!70] (4.6,0) circle (0.028);
\fill[black!70] (4.81,-0.21) circle (0.028);
\fill[black!70] (5.02,-0.42) circle (0.028);
\fill[black!70] (5.23,-0.21) circle (0.028);
\fill[black!70] (5.44,0) circle (0.028);
\fill[black!70] (5.65,0.21) circle (0.028);
\fill[black!70] (5.86,0) circle (0.028);
\fill[black!70] (6.07,0.21) circle (0.028);
\fill[cBad] (4.81,-0.21) circle (0.075);
\node[font=\small] at (5.335,-0.865) {$j=2$\ \textcolor{cBad}{$\times$}};
\draw[gray!60,dashed,line width=0.4pt] (6.82,0) -- (8.45,0);
\draw[cA,line width=1.2pt,line cap=round] (6.9,0) -- (7.11,-0.21);
\draw[gray!75!black,line width=1.2pt,line cap=round] (7.11,-0.21) -- (7.32,0);
\draw[gray!75!black,line width=1.2pt,line cap=round] (7.32,0) -- (7.53,0.21);
\draw[gray!75!black,line width=1.2pt,line cap=round] (7.53,0.21) -- (7.74,0.42);
\draw[cA,line width=1.2pt,line cap=round] (7.74,0.42) -- (7.95,0.21);
\draw[gray!75!black,line width=1.2pt,line cap=round] (7.95,0.21) -- (8.16,0.42);
\draw[cA,line width=1.2pt,line cap=round] (8.16,0.42) -- (8.37,0.21);
\fill[black!70] (6.9,0) circle (0.028);
\fill[black!70] (7.11,-0.21) circle (0.028);
\fill[black!70] (7.32,0) circle (0.028);
\fill[black!70] (7.53,0.21) circle (0.028);
\fill[black!70] (7.74,0.42) circle (0.028);
\fill[black!70] (7.95,0.21) circle (0.028);
\fill[black!70] (8.16,0.42) circle (0.028);
\fill[black!70] (8.37,0.21) circle (0.028);
\fill[cBad] (7.11,-0.21) circle (0.075);
\node[font=\small] at (7.635,-0.865) {$j=3$\ \textcolor{cBad}{$\times$}};
\fill[green!50!black!8] (8.98,-0.64) rectangle (10.89,0.756);
\draw[green!50!black,line width=0.7pt] (8.98,-0.64) rectangle (10.89,0.756);
\draw[gray!60,dashed,line width=0.4pt] (9.12,0) -- (10.75,0);
\draw[gray!75!black,line width=1.2pt,line cap=round] (9.2,0) -- (9.41,0.21);
\draw[gray!75!black,line width=1.2pt,line cap=round] (9.41,0.21) -- (9.62,0.42);
\draw[gray!75!black,line width=1.2pt,line cap=round] (9.62,0.42) -- (9.83,0.63);
\draw[cA,line width=1.2pt,line cap=round] (9.83,0.63) -- (10.04,0.42);
\draw[gray!75!black,line width=1.2pt,line cap=round] (10.04,0.42) -- (10.25,0.63);
\draw[cA,line width=1.2pt,line cap=round] (10.25,0.63) -- (10.46,0.42);
\draw[cA,line width=1.2pt,line cap=round] (10.46,0.42) -- (10.67,0.21);
\fill[black!70] (9.2,0) circle (0.028);
\fill[black!70] (9.41,0.21) circle (0.028);
\fill[black!70] (9.62,0.42) circle (0.028);
\fill[black!70] (9.83,0.63) circle (0.028);
\fill[black!70] (10.04,0.42) circle (0.028);
\fill[black!70] (10.25,0.63) circle (0.028);
\fill[black!70] (10.46,0.42) circle (0.028);
\fill[black!70] (10.67,0.21) circle (0.028);
\node[font=\small] at (9.935,-0.865) {$j=4$\ \textcolor{green!50!black}{\checkmark}};
\draw[gray!60,dashed,line width=0.4pt] (11.42,0) -- (13.05,0);
\draw[gray!75!black,line width=1.2pt,line cap=round] (11.5,0) -- (11.71,0.21);
\draw[gray!75!black,line width=1.2pt,line cap=round] (11.71,0.21) -- (11.92,0.42);
\draw[cA,line width=1.2pt,line cap=round] (11.92,0.42) -- (12.13,0.21);
\draw[gray!75!black,line width=1.2pt,line cap=round] (12.13,0.21) -- (12.34,0.42);
\draw[cA,line width=1.2pt,line cap=round] (12.34,0.42) -- (12.55,0.21);
\draw[cA,line width=1.2pt,line cap=round] (12.55,0.21) -- (12.76,0);
\draw[gray!75!black,line width=1.2pt,line cap=round] (12.76,0) -- (12.97,0.21);
\fill[black!70] (11.5,0) circle (0.028);
\fill[black!70] (11.71,0.21) circle (0.028);
\fill[black!70] (11.92,0.42) circle (0.028);
\fill[black!70] (12.13,0.21) circle (0.028);
\fill[black!70] (12.34,0.42) circle (0.028);
\fill[black!70] (12.55,0.21) circle (0.028);
\fill[black!70] (12.76,0) circle (0.028);
\fill[black!70] (12.97,0.21) circle (0.028);
\fill[cBad] (12.76,0) circle (0.075);
\node[font=\small] at (12.235,-0.865) {$j=5$\ \textcolor{cBad}{$\times$}};
\draw[gray!60,dashed,line width=0.4pt] (13.72,0) -- (15.35,0);
\draw[gray!75!black,line width=1.2pt,line cap=round] (13.8,0) -- (14.01,0.21);
\draw[cA,line width=1.2pt,line cap=round] (14.01,0.21) -- (14.22,0);
\draw[gray!75!black,line width=1.2pt,line cap=round] (14.22,0) -- (14.43,0.21);
\draw[cA,line width=1.2pt,line cap=round] (14.43,0.21) -- (14.64,0);
\draw[cA,line width=1.2pt,line cap=round] (14.64,0) -- (14.85,-0.21);
\draw[gray!75!black,line width=1.2pt,line cap=round] (14.85,-0.21) -- (15.06,0);
\draw[gray!75!black,line width=1.2pt,line cap=round] (15.06,0) -- (15.27,0.21);
\fill[black!70] (13.8,0) circle (0.028);
\fill[black!70] (14.01,0.21) circle (0.028);
\fill[black!70] (14.22,0) circle (0.028);
\fill[black!70] (14.43,0.21) circle (0.028);
\fill[black!70] (14.64,0) circle (0.028);
\fill[black!70] (14.85,-0.21) circle (0.028);
\fill[black!70] (15.06,0) circle (0.028);
\fill[black!70] (15.27,0.21) circle (0.028);
\fill[cBad] (14.22,0) circle (0.075);
\node[font=\small] at (14.535,-0.865) {$j=6$\ \textcolor{cBad}{$\times$}};
\end{tikzpicture}}

% =====================================================================
% 附錄 A 第 1 段：A.1–A.4（動畫 M01–M04）
% 圖由 mit/make_figs_part1.py 產生（含數字檢查），請勿手改 mit/figs_part1.tex
% =====================================================================

% ---------------------------------------------------------------------
\subsection{Cauchy 乘積}\label{sec:m01}

第 1.12 節把「根＋左＋右」直接翻譯成 $C=1+xC^2$。這一小節把中間的關鍵一步看清楚：
為什麼「左、右互相獨立」在代數上就是「生成函數相乘」\mitref{L1}。

\textbf{冪級數的乘法。}設 $A(x)=\sum_n a_nx^n$、$B(x)=\sum_n b_nx^n$。相乘時
$a_kx^k$ 和 $b_{n-k}x^{n-k}$ 的乘積正好落在 $x^n$，所以
\begin{equation}\label{eq:m01-cauchy}
[x^n]\,(A\cdot B)=\sum_{k=0}^{n}a_k\,b_{n-k}=a_0b_n+a_1b_{n-1}+\cdots+a_nb_0 .
\end{equation}
這叫做 \textbf{Cauchy 乘積}。排成乘法表最清楚：第 $i$ 列放 $a_ix^i$、第 $j$ 行放 $b_jx^j$，
格子 $(i,j)$ 放 $a_ib_j\,x^{i+j}$。次方相同的格子都在同一條反對角線 $i+j=n$ 上，
$[x^n]$ 就是把這條反對角線上的格子加起來。例如 $x^2$ 的係數（圖~\ref{fig:m01-table} 的黃色格子）是
$\textcolor{cLeft}{C_0}\textcolor{cRight}{C_2}+\textcolor{cLeft}{C_1}\textcolor{cRight}{C_1}
+\textcolor{cLeft}{C_2}\textcolor{cRight}{C_0}=2+1+2=5$。

\begin{figure}[H]
\centering
\FigMiCauchy
\caption{$C(x)\cdot C(x)$ 的乘法表：\textcolor{cLeft}{列}是第一個 $C(x)$ 的項，\textcolor{cRight}{行}是第二個。
同一條反對角線 $i+j=n$ 的格子同色，加起來就是 $[x^n]\,C^2$。淡灰色的格子屬於 $x^4$ 以上，
這張 $4\times4$ 的表只畫到它們的一部分。}
\label{fig:m01-table}
\end{figure}

\textbf{組合意義。}$\textcolor{cLeft}{C_k}\textcolor{cRight}{C_{n-k}}$ 數的是有序對
「\textcolor{cLeft}{左邊}一個大小 $k$ 的卡塔蘭物件、\textcolor{cRight}{右邊}一個大小 $n-k$ 的」：
左右各自選、互不影響，所以\textbf{相乘}；$k$ 不同是不同的情況，所以\textbf{相加}。
因此 $[x^n]\,C^2$ 就是「總大小為 $n$ 的（左，右）有序對」的個數。
拿掉根之後剩下的左、右兩部分，正好就是這樣的有序對：把遞迴式 \eqref{eq:segner} 的 $n$ 換成 $n+1$，
\begin{equation}\label{eq:m01-shift}
C_{n+1}=\sum_{k=0}^{n}\textcolor{cLeft}{C_k}\,\textcolor{cRight}{C_{n-k}}=[x^n]\,C^2\qquad(n\ge0).
\end{equation}
也就是說，$C^2$ 的係數整串往後移一格，就是 $C$ 去掉常數項：
\[
\begin{array}{r@{\;}c@{\;}*{5}{r@{\;}c@{\;}}l}
C^2 &=& 1&+&2x&+&5x^2&+&14x^3&+& & &\cdots\\
xC^2 &=& & & x&+&2x^2&+&5x^3&+&14x^4&+&\cdots\\
C &=& 1&+&x&+&2x^2&+&5x^3&+&14x^4&+&\cdots
\end{array}
\]
所以 $C-1=xC^2$，也就是 $C=1+xC^2$，即式 \eqref{eq:cgf}。

\begin{keybox}[獨立 $\to$ 相乘]
兩個互相獨立的部分 $\longleftrightarrow$ 生成函數相乘（Cauchy 乘積 \eqref{eq:m01-cauchy}）；
根本身佔大小 $1$ $\longleftrightarrow$ 乘上 $x$；空的物件 $\longleftrightarrow$ 常數 $1$。
這是第 5 節「兩條翻譯規則」只有一個變數的情形。
\end{keybox}

% ---------------------------------------------------------------------
\subsection{生成函數的閉式}\label{sec:m02}

第 1.12 節解出 $C(x)$ 並直接列出結果。這裡補上兩個細節 \mitref{L1}：
\textbf{為什麼取減號}，以及 \textbf{$\sqrt{1-4x}$ 的係數怎麼算}。
$C=1+xC^2$ 就是 $xC^2-C+1=0$，由二次公式（$a=x$，$b=-1$，$c=1$）得 $C=\bigl(1\pm\sqrt{1-4x}\,\bigr)/(2x)$。

\noindent
\begin{minipage}[c]{0.56\textwidth}
\textbf{取哪一個號？}冪級數 $C(x)=C_0+C_1x+\cdots$ 在 $x=0$ 的值是常數項 $C_0=1$。
\begin{itemize}[leftmargin=1.2em]
  \item 取 $+$：$x\to0$ 時分子趨近 $2$，所以
        $\frac{1+\sqrt{1-4x}}{2x}\approx\frac{2}{2x}=\frac1x\to\infty$，
        不可能是 $C_0+C_1x+\cdots$。
  \item 取 $-$：用 $\sqrt{1-4x}\approx1-2x$，得 $\frac{1-(1-2x)}{2x}=1=C_0$。
        也可以先有理化：$\frac{1-\sqrt{1-4x}}{2x}=\frac{2}{1+\sqrt{1-4x}}$，代 $x=0$ 就是 $1$。
\end{itemize}
\end{minipage}\hfill
\begin{minipage}[c]{0.4\textwidth}
\centering
\FigMiiBranches
\captionof{figure}{兩個根在 $0<x\le\frac14$ 的圖形；兩支在 $x=\frac14$（根號為 $0$）會合於 $2$。}
\label{fig:m02-branches}
\end{minipage}

\medskip
\textbf{二項級數。}Newton 的廣義二項式定理：$(1+u)^r=\sum_{m\ge0}\binom{r}{m}u^m$，
其中 $\binom{r}{m}=\frac{r(r-1)\cdots(r-m+1)}{m!}$，$r$ 可以是任意實數。
取 $r=\frac12$、$u=-4x$，得到 $\sqrt{1-4x}=\sum_{m\ge0}\binom{1/2}{m}(-4)^m\,x^m$。前幾項：
\begin{center}
\small
\renewcommand{\arraystretch}{1.2}
\begin{tabular}{@{}c|cccccc@{}}
$m$ & 0 & 1 & 2 & 3 & 4 & 5\\
\hline
$\binom{1/2}{m}$ & $1$ & $\frac12$ & $-\frac18$ & $\frac1{16}$ & $-\frac5{128}$ & $\frac7{256}$\\
$(-4)^m$ & $1$ & $-4$ & $16$ & $-64$ & $256$ & $-1024$\\
\hline
$\binom{1/2}{m}(-4)^m$ & $\mathbf{1}$ & $\mathbf{-2}$ & $\mathbf{-2}$ & $\mathbf{-4}$ & $\mathbf{-10}$ & $\mathbf{-28}$
\end{tabular}
\end{center}
分數乘上 $(-4)^m$ 之後全都是整數。\textbf{一般項}（$m\ge1$）：分子的 $m$ 個因子
$\frac12,-\frac12,-\frac32,\ldots,-\frac{2m-3}{2}$ 各帶一個 $\frac12$，其中 $m-1$ 個是負的，
所以 $\binom{1/2}{m}=\frac{(-1)^{m-1}\,1\cdot3\cdots(2m-3)}{2^m\,m!}$；乘上 $(-4)^m$，
再把奇數連乘積的分子、分母同乘偶數 $2\cdot4\cdots(2m-2)$：$1\cdot3\cdots(2m-3)=\frac{(2m-2)!}{2\cdot4\cdots(2m-2)}=\frac{(2m-2)!}{2^{m-1}(m-1)!}$
（$m=1$ 時是空乘積 $1$）：
\begin{equation}\label{eq:m02-coef}
\binom{1/2}{m}(-4)^m=-\frac{2^m\cdot1\cdot3\cdots(2m-3)}{m!}=-\frac{2\,(2m-2)!}{m!\,(m-1)!}
=-\frac{2}{m}\binom{2m-2}{m-1}\qquad(m\ge1).
\end{equation}

\textbf{代回去。}
\begin{align*}
\sqrt{1-4x}&=1-2x-2x^2-4x^3-10x^4-28x^5-\cdots,\\
C(x)&=\frac{1-\sqrt{1-4x}}{2x}=\frac{2x+2x^2+4x^3+10x^4+28x^5+\cdots}{2x}=1+x+2x^2+5x^3+14x^4+\cdots
\end{align*}
一般地，$C_n$ 是 $-\frac12$ 乘上 $\sqrt{1-4x}$ 的 $x^{n+1}$ 係數；在式 \eqref{eq:m02-coef} 取 $m=n+1$：
\[
\boxed{\;C_n=-\frac12\,[x^{n+1}]\sqrt{1-4x}=-\frac12\cdot\Bigl(-\frac{2}{n+1}\binom{2n}{n}\Bigr)=\frac{1}{n+1}\binom{2n}{n}\;}
\]

% ---------------------------------------------------------------------
\subsection{醉漢走路}\label{sec:m03}

醉漢站在數線的 $0$，每一步往右 $+1$ 或往左 $-1$，機率各 $\frac12$，各步互相獨立。
$-1$ 是懸崖：走到 $-1$ 就掉下去。\textbf{他最後一定會掉下去嗎？} \mitref{L1--L2}
把時間畫成橫軸、位置畫成縱軸，往上一步記 U、往下一步記 D，一次走法就是一條 $\pm1$ 路徑。

\textbf{為什麼是 Dyck 路徑？}假設第 $t$ 步第一次走到 $-1$，那麼之前一直 $\ge0$；一步只走 $1$，
所以第 $t-1$ 步一定在 $0$、第 $t$ 步往下。從 $0$ 回到 $0$ 要走偶數步，所以 $t-1=2n$；
前 $2n$ 步「一直 $\ge0$、最後回到 $0$」正是 Dyck 路徑。反過來，半長 $n$ 的 Dyck 路徑再往下一步，
就是在第 $2n+1$ 步第一次掉下去。

\begin{thmbox}[第一次掉下去的時刻]
第一次走到 $-1$ 發生在第 $2n+1$ 步 $\iff$ 前 $2n$ 步是半長 $n$ 的 Dyck 路徑，而第 $2n+1$ 步往下。
這樣的走法有 $C_n$ 條，每一條的機率都是 $(\frac12)^{2n+1}$，所以
\[
P(\text{第 $2n+1$ 步第一次掉下去})=\frac{C_n}{2^{2n+1}} .
\]
\end{thmbox}

$n=0,1,2,3$ 時機率分別是 $\frac12,\frac18,\frac2{32},\frac5{128}$（圖~\ref{fig:m03-walk}）。

\begin{figure}[H]
\centering
\FigMiiiWalk
\caption{左：走法 \textcolor{cA}{UUDUDD}\,\textcolor{cBad}{D}（貫穿例子的 Dyck 路徑再往下一步），第 7 步第一次走到 $-1$。
右：第 $2n+1$ 步第一次掉下去的所有走法（$n=0,1,2,3$），\textcolor{cA}{框起來}的就是左邊那一條；
灰線是地面 $0$，\textcolor{cBad}{紅虛線}是懸崖 $-1$。}
\label{fig:m03-walk}
\end{figure}

\textbf{全部加起來。}不同的 $n$ 是互斥的情況，用上一小節的閉式：
\[
P(\text{最後掉下去})=\sum_{n\ge0}\frac{C_n}{2^{2n+1}}=\frac12\sum_{n\ge0}C_n\Bigl(\frac14\Bigr)^n=\frac12\,C\Bigl(\frac14\Bigr)
=\frac12\cdot\frac{1-\sqrt{1-1}}{2\cdot\frac14}=\frac12\cdot2=1 .
\]
$x=\frac14$ 正好是根號變成 $0$ 的地方，也就是收斂半徑的邊界。因為 $C_n\approx\frac{4^n}{\sqrt{\pi}\,n^{3/2}}$，
級數 $\sum C_n/4^n$ 仍然收斂，可以直接代入（Abel 定理）。也可以不用 Abel 定理，直接看部分和：

\begin{notebox}[部分和]
令 $S_K$ 為前 $2K-1$ 步之內掉下去的機率（$n<K$ 的項相加）。用歸納法可以驗證
\[
1-S_K=\frac{1}{4^K}\binom{2K}{K}=\frac12\cdot\frac34\cdot\frac56\cdots\frac{2K-1}{2K}\le\frac{1}{\sqrt{2K+1}}\longrightarrow0 .
\]
歸納的一步是 $\frac{1}{4^{K}}\binom{2K}{K}-\frac{1}{4^{K+1}}\binom{2K+2}{K+1}=\frac{C_K}{2^{2K+1}}$；
不等式來自 $\bigl(\frac{2i-1}{2i}\bigr)^2\le\frac{2i-1}{2i}\cdot\frac{2i}{2i+1}$。所以總和是 $1$，只是趨近得很慢：
$S_1,S_2,S_3,S_4=\frac12,\frac58,\frac{11}{16},\frac{93}{128}$（還沒掉下去的是 $\frac12,\frac38,\frac5{16},\frac{35}{128}$），
走了 199 步（$K=100$）還有約 $5.6\%$ 的機率沒掉下去。醉漢一定會掉下懸崖，只是可能要走很久。
\end{notebox}

% ---------------------------------------------------------------------
\subsection{循環移位}\label{sec:m04}

第 1.6 節用反射原理證明 $C_n=\binom{2n}{n}-\binom{2n}{n-1}$。這裡介紹另一個證明——
\textbf{循環移位} \mitref{L2}，它直接得到 $C_n=\frac{1}{2n+1}\binom{2n+1}{n}$，不需要做減法。

取 $n+1$ 個 $+1$ 與 $n$ 個 $-1$（總和 $+1$、長度 $2n+1$），首尾相接排成一圈。
從第 $j$ 點出發讀一圈（$j=0,1,\ldots,2n$），得到 $2n+1$ 種\textbf{循環讀法}；
如果出發之後每一個前綴和都 $>0$，就說這種讀法是\textbf{好的}。

\begin{thmbox}[古典循環引理]
$n+1$ 個 $+1$、$n$ 個 $-1$ 排成一圈，$2n+1$ 種讀法中\textbf{恰有一種}是好的。
\end{thmbox}

\textbf{例。}$n=3$，取下面第一列的序列；第二列是它從第 0 點到第 7 點的高度（前綴和）：
\[
\begin{array}{r@{\qquad}l}
\text{序列（$j=0$）} & -1,+1,-1,-1,+1,+1,+1\\
\text{高度} & 0,-1,0,-1,-2,-1,0,1\\
\text{從第 4 點讀起（$j=4$）} & +1,+1,+1,-1,+1,-1,-1
\end{array}
\]
最低點是第 4 點（高度 $-2$）。7 種讀法中只有 $j=4$ 是好的（圖~\ref{fig:m04-rotations}），
它是一個 $+1$ 接上貫穿例子的 Dyck 路徑 UUDUDD。

\begin{figure}[H]
\centering
\FigMivRotations
\caption{序列 $-1,+1,-1,-1,+1,+1,+1$ 的 7 種循環讀法（從第 $j$ 點出發）。灰色是 $+1$，\textcolor{cA}{藍色}是 $-1$；
\textcolor{cBad}{紅點}是第一個高度 $\le0$ 的點。只有 $j=4$ 是好的，它的起點正是原序列（$j=0$）的
\textcolor{cTop}{最後一個最低點}（高度 $-2$）。}
\label{fig:m04-rotations}
\end{figure}

\textbf{證明。}設各點高度是 $S_0=0,S_1,\ldots,S_{2n+1}=1$，令 $i$ 是 $S_0,\ldots,S_{2n}$ 中\textbf{最後一個最低點}，
最低的高度是 $\mu$。
\begin{itemize}
  \item \textbf{從 $i$ 出發是好的。}讀到原序列的第 $k$ 點（$i<k\le2n+1$）時，前綴和是 $S_k-S_i>0$
        （$i$ 之後沒有一樣低或更低的點，而 $S_{2n+1}=1>0\ge\mu$）；
        繞過接頭之後，讀到第 $k$ 點（$0<k\le i$）時前綴和是 $1+S_k-\mu\ge1$。
  \item \textbf{只有 $i$。}若從 $j$ 出發是好的，則 $k>j$ 時 $S_k>S_j$；$k<j$ 時 $1+S_k-S_j>0$，即 $S_k\ge S_j$。
        所以 $S_j$ 是最低的高度，而且 $j$ 之後沒有一樣低的點：$j=i$。\hfill$\square$
\end{itemize}

\textbf{計數。}這樣的序列共有 $\binom{2n+1}{n}$ 種（選出 $-1$ 的位置）。同一圈的 $2n+1$ 種讀法都不相同
（若有兩種相同，序列就有週期，好的讀法會不只一種），所以序列恰好分成 $\binom{2n+1}{n}/(2n+1)$ 圈，每圈恰有一種好的。
另一方面，好的序列第一步一定是 $+1$，之後一直 $\ge1$、最後回到 $1$：正是\textbf{一個 $+1$ 接一條 Dyck 路徑}，
共 $C_n$ 種。所以
\begin{equation}\label{eq:m04-cycle}
\boxed{\;C_n=\frac{1}{2n+1}\binom{2n+1}{n}=\frac{(2n)!}{n!\,(n+1)!}=\frac{1}{n+1}\binom{2n}{n}\;}
\end{equation}
例如 $n=3$：$\binom{7}{3}=35$ 種序列，每圈 7 種，恰好 $35/7=5=C_3$ 圈（另外 30 種是不好的）。
好的序列也就是第 1.9 節「A 得 $n+1$ 票、B 得 $n$ 票、A 始終嚴格領先」的開票順序。

\textbf{推廣。}第~\ref{sec:raney} 節的循環引理把同樣的想法推廣到終點高度 $t$ 與多種步長（$-1,-2,\ldots$）：
$N+t$ 種讀法中恰有 $t$ 種是好的，由此證明 Raney 公式。取 $\vv=(2)$、$t=1$、$\vn=(n)$ 就回到這裡：
$A_1(\vn)=\frac{1}{2n+1}\binom{2n+1}{n}=C_n$。
